\section{Decisões tomadas e por quê}

Nesta etapa, consolidamos as seguintes decisões técnicas e organizacionais:

\paragraph{1. Seleção do dataset oficial: Microdados do ENEM + gerador sintético.}
Optou-se pelos Microdados do ENEM (INEP) como conjunto de dados real oficial. A
escolha traz significado prático às métricas, permite trabalhar com um volume
expressivo de dados reais nacionais ($\sim 2$\,GB descompactados por ano) e
possibilita uma estratégia híbrida: o processamento real demonstrará a
utilidade prática do sistema, enquanto o gerador sintético (com mesma estrutura
de colunas e \emph{seed} fixa) será utilizado nos testes formais de \emph{speedup}
e escalabilidade da Lei de Amdahl.

\paragraph{2. Pergunta que o sistema responde.}
Definiu-se formalmente a pergunta analítica:
\begin{quote}
  \emph{``Qual o desempenho médio nas provas objetivas (Matemática) e de Redação
  dos participantes do exame, discriminado por Unidade Federativa (UF) e por
  dependência administrativa da escola (pública versus privada)?''}
\end{quote}
\textbf{Justificativa técnica para Sistemas Distribuídos:}
A pergunta selecionada exige filtragem (eliminar faltosos e eliminados) e
agregação por chave composta $(\text{UF}, \text{tipo de escola})$. No modelo
MapReduce, essa operação é altamente comutativa e associativa: cada worker
calcula as tuplas locais $(\text{soma\_notas}, \text{quantidade})$ e devolve ao
mestre uma estrutura de apenas $27 \times 2 = 54$ entradas. Isso garante que a
mensagem de retorno tenha tamanho constante de poucos kilobytes, preservando o
canal de rede e evitando que a fase de agregação final se torne um gargalo
prematuro.

\paragraph{3. Linguagem e biblioteca de comunicação.}
Decidiu-se pela utilização de **Python 3 com a biblioteca \texttt{mpi4py}** sobre
o Open MPI. Embora C ofereça menor consumo de memória por nó, o uso de Python
acelera drasticamente o desenvolvimento do código de ingestão, manipulação de
dicionários e validação dos dados, mantendo excelente desempenho através de
chamadas nativas do MPI subjacente.

\paragraph{4. Atribuição dos papéis fixos da equipe.}
Para assegurar continuidade e responsabilidade direta em cada frente ao longo
das semanas, foram distribuídos os papéis fixos:
\begin{itemize}
  \item \textbf{Dados e Documentação:} Matheus Nardi — responsável pelas
    rotinas de ingestão, sanitização, verificação de dados, consolidação dos
    relatórios e slides semanais;
  \item \textbf{Desenvolvimento:} Alêkson Castro — responsável pela
    implementação do algoritmo sequencial de referência, núcleo MPI (mestre e
    workers), fila dinâmica de tarefas e tolerância a falhas;
  \item \textbf{Infraestrutura e Experimentos:} Italo Beckman — responsável pelo
    provisionamento via Vagrant e Ansible, orquestração do cluster de VMs,
    configuração de rede dos nós físicos e coleta sistemática de métricas de
    desempenho.
\end{itemize}
